\section{Canonical refill and the conditional arithmetic point}
\label{sec:canonical-refill}

\status{RF1--RF8 construct the literal canonical refill, complete the
\texttt{exists\_ledPart} implication using LP, and install the resulting
ideal-scaled lattice and height in the point. RF9 proves the two scale
comparisons needed by that construction. RF10 proves the uniform point join
from the explicitly retained recursive, host-column, level and initial-size
suppliers. The geometric \texttt{Lay}/\texttt{RecOK}/\texttt{LedOK} package,
the analytic supply of level primes and the stronger \texttt{NumAt} record
are not inferred from the refill. RF11--RF15 additionally prove the
quantitative rank, thin-size, gap and band estimates and construct every
field of the block-length record \texttt{ParamsT}; RF16 proves the
actual upper bound on word cardinality. The complete
\texttt{ParameterChoice'} and \texttt{SizesAt} suppliers remain explicit.}

Fix the actual datum $D$, a size datum $L_0$, the local constants $K$,
and an error datum $\mathcal E$. Write $p_w$ and $e_w>0$ for the residue
characteristic and ramification index of a finite place. The set of words
$u$ is finite. Here $m=D.R.m$ is the global chain length, $k=D.R.k$,
$n=|A|$, and $J=\{1,\ldots,m\}$. Put
\[
 b=e^{x^\epsilon},\quad a=xb,\quad X'=e^{x^{\epsilon/3}},\quad
 b'=a/X',\quad w''=b'/(xn),\quad X''=a/w''=X'xn.       \tag{RF.1}
\]
The top degree functions $a_T,b_T$, the weights $w_j(\tau)$, and all
geometric and representation data are the supplied ones. In particular,
this construction will retain $L_0$'s scale and its exact selected degrees.

\paragraph{RF1: choose the same representative throughout an orbit.}
The finite abelian group $A$ comes with its source enumeration
$\operatorname{ofNat}$ and integer labels. For $h,t\in A$, define
$\operatorname{bs}(h,t)$ to be $\operatorname{ofNat}(v)$ for the least
nonnegative $v$ such that
\[
                  \operatorname{ofNat}(v)+ih=t
                  \quad\hbox{for some }i\geq0.                 \tag{RF.2}
\]
Such a $v$ exists by taking the label of $t$. The same candidate labels
work for $t$ and $t+h$: add one to $i$ in one direction, and add
$\operatorname{ord}(h)-1$ in the other. Thus
\[
 \operatorname{bs}(h,t)+ih=t\ \hbox{for some }i\geq0,
 \qquad\operatorname{bs}(h,t+h)=\operatorname{bs}(h,t).          \tag{RF.3}
\]
This is the literal \texttt{Arith.bsH}; its choice uses the given labels,
so it remains compatible when points are compared in those labels.

\paragraph{RF2: fill the actual weights, including the unused indices.}
Label the archimedean positions by $z=(\delta,t)\in\{0,1\}\times A$,
using the original bijection $\tau(z)$. A top is $(0,t)$ with
$t\in TT$; a thick position has first coordinate $1$ and second
coordinate in the translated second set; all other positions are thin.
These are the actual top and thick sets: injectivity of the position map
shows that their label descriptions agree with the geometric descriptions.
Extend $b_T$ by zero away from $TT$ and set
\[
 o_j(z)=
 \begin{cases}
 w_j(\tau(z))-\mathbf1_{\mathrm{top}}(z)b_T(t),&j\in J,\\
 b'/x,&j\notin J,\ z\text{ top},\\
 b,&j\notin J,\ z\text{ thick},\\
 w'',&j\notin J,\ z\text{ thin}.
 \end{cases}                                                  \tag{RF.4}
\]
The three supplied weight windows give, respectively,
\[
 \frac{b'}x\leq o_j(z)\leq\frac{b'}x+\eta,\qquad
 b\leq o_j(z)\leq b+\eta+1,\qquad
 w''\leq o_j(z)\leq w''+\eta.                                 \tag{RF.5}
\]
They apply for every $j$, because (RF.4) gives their lower endpoints
outside $J$. Define
\[
 \operatorname{fill}_j(z)=\mathbf1_{\mathrm{top}}(z)b_T(t)+o_j(z).
\]
Then $\operatorname{fill}_j(z)=w_j(\tau(z))$ for $j\in J$.
If $w''\geq1$, every fill is at least $1$: thin positions use $w''$,
thick positions use $b\geq1$, and tops use $b_T\geq b'>0$ and
$b'/x=nw''\geq1$.

Use (RF.3) for the orbit selector, (RF.4) for its offsets, and set
\[
 e_p=\left\lceil C_e f k^2\log_p a\right\rceil_+,
 \qquad \kappa_I=\inf\{\kappa:\kappa\text{ is an admissible
                            universal Type-I loss}\}.         \tag{RF.6}
\]
The admissible set is the actual \texttt{kapISet}; its nonemptiness and
infimum attainment were proved in GJ1. Defining (RF.6) gives
\texttt{EpDef} and \texttt{KapIDef}; it does not by itself identify $e_p$
with the geometric stabilizer sandwich exponent. That distinct assertion
is \texttt{EpTie}. Equations (RF.3)--(RF.4) prove \texttt{CanonInst},
including its prescribed first-coordinate-zero offsets outside $J$.

\paragraph{RF3: one finite set supports every numerical field.}
Assume that the twist set
$S_{\rm tw}=\{w:\mu_w\ne\mu_0\}$ is finite and that the positive
integers $D_{\Gamma,u}$ are supplied. Let $T_u$ denote the actual type
prime of word $u$, and choose
\[
 B_S=\lceil X''\rceil_++g_0+g_1+d+
                           \sum_u(T_u+D_{\Gamma,u}),\qquad
 S=\{w:p_w\leq B_S\}\cup S_{\rm tw}.                          \tag{RF.7}
\]
This set is finite. Indeed, a place over a rational prime $p$ determines
a prime ideal dividing the nonzero principal ideal $(p)$ in the number
ring; unique ideal factorization makes that set finite. There are only
finitely many $p\leq B_S$.
If $w\notin S$, then
\[
 \begin{gathered}
 p_w>X'',\quad p_w\notin\{g_0,g_1,d\},\quad
 p_w\ne T_u,\quad p_w\nmid D_{\Gamma,u}\quad(\text{every }u),\\
                         \mu_w=\mu_0.                         \tag{RF.8}
 \end{gathered}
For the divisibility assertion use $D_{\Gamma,u}>0$: a positive integer
divisible by $p_w$ is at least $p_w$, contradicting $p_w>B_S$.
In particular $S$ contains every place of residue characteristic at most
$X''$, as well as every nontrivially twisted place.

\paragraph{RF4: the literal gain, forcing, error and deficit values.}
We spell out the outer cases, since they are what makes the refill
possible. Let $\mathcal R_p$ be the orbit representatives from RF1 and
$\mathcal V_p=\texttt{VV}(p)$ the actual served representatives. Write
$g^0_{jv},g^U_{jv},P_{jv}$ for the source local gain scores and correction,
$t_v,D_p$ for the local depth and degree budgets, and
$\Sigma_p=m\sum_{v\in\mathrm{goodReps}(p)}\kappa_v$.
All are evaluated at the same selected $a_T,b_T$, scale and instance.
The explicit formulas for these unchanged scalar scores are recorded in
the expanded report. Set
\[
 S_v=\left\lceil\sum_{j\in J}
        \max\{0,g^U_{jv}-P_{jv}-\max(g^0_{jv},0)\}\right\rceil,
                                                               \tag{RF.9}
\]
and define
\[
 \begin{split}
 \tau_p&=\begin{cases}0,&\mathcal V_p=\varnothing,\\
 e_p+m\sum_{v\in\mathcal V_p}(\lceil t_v\rceil_++D_p+1),
                                      &\mathcal V_p\ne\varnothing,
                         \end{cases}\\
 \iota_p&=\begin{cases}0,&\mathcal V_p=\varnothing,\\
 2\Sigma_p\max_{v\in\mathcal V_p}S_v,&\mathcal V_p\ne\varnothing.
                         \end{cases}                          \tag{RF.10}
 \end{split}
\]
The gain is zero when $p>X''$ or when $p$ is of the source inert class
and $h_p=0$. Otherwise it is
\[
 G_{jp}=\sum_{v\in\mathcal V_p}
       \max\{g^U_{jv}-P_{jv},\max(g^0_{jv},0)\}
          +\sum_{v\in\mathcal R_p\setminus\mathcal V_p}
                                      \max(g^0_{jv},0).         \tag{RF.11}
\]
At a place $w$, install $G_{j,p_w}$ only for the inert or signal class,
and install zero for the other classes. In particular it vanishes at
$g_0,g_1$: \texttt{HostArith} places these level primes in the split class.

The forcing value is, in priority order,
\[
 F_{jp}=\begin{cases}
 F^G_j,&p\in\{g_0,g_1\},\\
 F^d_j,&\operatorname{cls}(p)=\mathrm{ramD},\\
 F^{\rm in}_{jp},&p\leq X'',\ p\text{ inert or signal},\\
 0,&\text{otherwise}.
 \end{cases}                                                  \tag{RF.12}
\]
Here $F^G_j=-\operatorname{fockDeg}_j/2$,
$F^d_j=-(k-1)\sum_t\min(w_j(0,t),w_j(1,t))/2$, and
$F^{\rm in}_{jp}$ is the explicit finite orbit/linear-program expression
in the report. In the error formula the own type prime has first priority:
\[
 E_{ujw}=\begin{cases}
 E^T_{uw},&p_w=T_u,\\
 v_{p_w}(D_{\Gamma,u})+E^G_{0,ujw},&p_w\ne T_u,\ p_w=g_0,\\
 v_{p_w}(D_{\Gamma,u})+E^G_{1,ujw},&p_w\ne T_u,g_0,\ p_w=g_1,\\
 v_{p_w}(D_{\Gamma,u})+E^d_j,&\text{next, if }\operatorname{cls}(p_w)=\mathrm{ramD},\\
 v_{p_w}(D_{\Gamma,u})+E^C_{jp_w},&\text{next, if }p_w\leq X'',\ p_w\text{ inert},\ h_{p_w}=0,\\
 v_{p_w}(D_{\Gamma,u})+E^I_{jp_w},&\text{next, if }p_w\leq X'',\ p_w\text{ inert or signal},\\
 v_{p_w}(D_{\Gamma,u}),&\text{otherwise}.
 \end{cases}                                                  \tag{RF.13}
\]
The report records every term of $E^T,E^G,E^d,E^C,E^I$, including the
ramification division in $E^G$ and the odd-order correction in $E^I$.
No sign or size estimate for these quantities is asserted merely by
installing their values.

Every field in (RF.10)--(RF.13) is zero off $S$. For $F$ and $E$, use
(RF.8), the implication $\operatorname{cls}(p)=\mathrm{ramD}\Rightarrow
p=d$, and $v_p(D_{\Gamma,u})=0$ when $p\nmid D_{\Gamma,u}$.
For $G$ use its $p>X''$ cutoff. A served prime must satisfy $p<X'$;
since $X'\leq X''$ by (RF.1), no place outside $S$ is served, giving
$\tau=\iota=0$ there. Conjugate places have the same rational residue
characteristic, so (RF.12) also gives $F_{j,\bar w}=F_{jw}$.

\paragraph{RF5: refill the record without changing its parameters.}
Replace the five numerical fields of $L_0$ by RF4 and its support by
$S$. Define
\[
 \begin{split}
 \mathrm{drop}&=\{p:\operatorname{NoGain}(p)\text{ and }p\ne T_u
                                             \text{ for every }u\},\\
 TT_p&=\begin{cases}TT,&p\text{ signal or }p\geq X',\\
 \{t\in TT:\operatorname{bs}(\operatorname{step}_p,t)\in\mathcal V_p\},
                                                       &\text{otherwise},
 \end{cases}\\
 w^{\rm part}_{j,t,p,e}&=\operatorname{fill}_j
             \bigl(\operatorname{partner}(p,t,e)\bigr).        \tag{RF.14}
 \end{split}
\]
Store $TT_p$ as the corresponding subset of the subtype $TT$. This
subtype introduces no lost labels because $TT_p\subset TT$; forgetting
the subtype gives exactly $TT_p$. Its full-set clause follows from the
first branch in (RF.14), and $w^{\rm part}\geq1$ follows from RF2.
The set $\mathrm{drop}$ is generally infinite. Finite support is asserted
for the five numerical fields, not for that set.

The remaining essential supplied condition is
\[
 \mu_w=\mu_0\quad\Longrightarrow\quad F_{jw}=0\quad(j\in J).
                                                               \tag{RF.15}
\]
Together with $L_0$'s untwisted vanishing of $A_j(w)$, this proves the
refilled untwisted field. At a place of neither inert nor signal class,
the gain is zero by definition and $\mathcal V_p=\varnothing$ by the
candidate condition; hence both deficits vanish. Type primes are excluded
from $\mathrm{drop}$ by definition.

All other fields are copied. The formulas used in RF1--RF4 depend on
the copied fields and on the fixed $D$, never on the fields being replaced.
Reconstructing the canonical instance over the refilled datum therefore
gives exactly the same selector, offsets, $\kappa_I$ and $e_p$.
The weight identity in RF2, partner identity and subset identity in
(RF.14), and the definitions of gain and deficits prove every clause of
\texttt{ParametersFill}. Its final clause $\mathrm{La}(K,s)=L.\mathrm{LaG}$
is the retained \texttt{Val.La\_eq}. Thus we obtain the literal
\texttt{exists\_refill} conclusion, preserving
\[
           L.sc=L_0.sc,\qquad L.a_T=L_0.a_T,\qquad L.b_T=L_0.b_T,
                                                               \tag{RF.16}
\]
and \texttt{Val}, together with \texttt{CanonInst}, \texttt{EpDef},
\texttt{KapIDef}, \texttt{FDef}, \texttt{EDef}, both support containments,
all-place \texttt{ParametersFill}, and the exact drop-set equivalence.

\paragraph{RF6: a served prime is never discarded.}
Nonempty $\mathcal V_p$ forces the literal candidate conditions
\[
 p\text{ prime},\quad p<X',\quad p\text{ inert},\quad
 x^{\epsilon/4}\leq\operatorname{ord}(h_p),\quad
 p\geq p_0\Longrightarrow p\notin\mathrm{Band},\quad
                         \Sigma_p\leq x^{-\epsilon}.           \tag{RF.17}
\]
The band's power condition is
$|p^i-X'|\leq\rho_B X'$ for some $i\geq1$, with $p_0\leq p<X'$.
Consequently (RF.17) excludes the small-order, power-band and excess-loss
alternatives in \texttt{NoGain}, as well as its $p>X''$ and
non-inert/non-signal alternatives. The only remaining alternative is
\texttt{IsCs}: inert class with nonzero odd-order $h_p$.

Here is the needed group argument. The frame has a homomorphism from
$A$ into a $2$-group, injective on a cyclic subgroup $A_0$, and
$[A:A_0]=o\leq x^{\epsilon/100}$. An odd-order element has zero image
in the $2$-group. Its multiple by $o$ lies in $A_0$ and still has zero
image; injectivity gives $oh_p=0$. Thus its order divides $o$, contrary to
\[
 x^{\epsilon/4}\leq\operatorname{ord}(h_p)\leq o
                 \leq x^{\epsilon/100}<x^{\epsilon/4}.
\]
We have proved
\[
 \mathcal V_p\ne\varnothing\ \Longrightarrow\
 \neg\operatorname{NoGain}(p),\qquad
 p\in\mathrm{drop}\ \Longrightarrow\ \mathcal V_p=\varnothing.
                                                               \tag{RF.18}
\]
This uses the actual small odd-part bound; a numerical lower bound on the
order alone would not imply even order in an arbitrary finite group.

\paragraph{RF7: complete the supplied-data \texttt{LedPart} construction.}
Assume the literal \texttt{Lay}, \texttt{RecOK} and \texttt{LedOK} data,
\texttt{Std}, $c_g\geq10$, an initial $L_0$ with the required
\texttt{Val} and degree identities, \texttt{HostArith}, and $w''\geq1$.
Retain also the source's $X'\leq T_u$ assumption. The relevant content of
\texttt{RecOK} is finite twist support, positive $D_{\Gamma,u}$, fullness
of each actual $M_{\rm ram}(u)$, and the two-sided commensurability at
its own type prime. The quantified \texttt{LedOK} supplies (RF.15) and,
for every canonical instance satisfying the defining numerical and fill
identities, the actual \texttt{LatTied}, \texttt{EpTie} and served-place
containment/commensurability predicates of LP4.

Apply \texttt{LedOK} first to $L_0$ to get (RF.15), then apply RF5.
Because (RF.16) and \texttt{Val} survive, apply \texttt{LedOK} again to
this refilled datum. This second application is essential: one cannot
transfer the geometric local predicates from an unrelated instance.
LP4 now patches the prescribed served and own-type local modules into
one $N_u$ for every word. LP5 installs these modules and
\[
 \operatorname{lenMN}(u,w)
       =\frac{\operatorname{toNat}(\operatorname{lenLoc}_w
                          (N_u,M_{\rm ram}(u)))}{e_w}.          \tag{RF.19}
\]
At a non-type place where $p\in\mathrm{drop}$, (RF.18) makes the place
unserved, so LP gives $N_{u,w}=M_{{\rm ram},u,w}$. The same is true
if $p$ is of neither inert nor signal class. At every non-type place
LP gives the required local containment. These are the first two clauses
of \texttt{NuDef}. Its third clause is the exact finite-length equation
$\ell=e_w\operatorname{lenMN}(u,w)$ from (RF.19).
Equality of localized modules makes their localized quotient zero, so
every non-type unserved length is zero. This proves \texttt{TiesAt.lats0}
and \texttt{NuUnserved}, not just the patch's existence.

The served and type identities are precisely \texttt{LatDefs}; all-place
filling and both support containments come from RF5; \texttt{LatTied} and
\texttt{EpTie} come from the second \texttt{LedOK} application. The
supplied level/error data and (RF.13) give \texttt{ErrDefs}; its three
constant inequalities follow from the standard fixed constants and
$c_g\geq10$. We have therefore constructed the entire literal
\texttt{LedPart}, with the same scale, selected degrees, \texttt{Val}
and canonical instance. No global containment $N_u\subset M_{\rm ram}(u)$
is added at a type prime, and finiteness of an arbitrary local length is
not asserted by (RF.19).

\paragraph{RF8: install the actual final lattice and height.}
For the constructed $L$ put
\[
 s_{uw}=\sum_{j\in J}\bigl(G^{\rm kept}_{jw}+F_{jw}-E_{ujw}\bigr),
 \qquad
 I_u=\prod_{w\in S}\mathfrak p_w^{-\lfloor e_ws_{uw}\rfloor},
 \qquad N'_u=I_uN_u.                                         \tag{RF.20}
\]
Here $G^{\rm kept}_{jw}=0$ on the drop set and is the installed gain
otherwise. The fractional-ideal product is finite by RF3--RF5; its powers
are integral and may have either sign. In particular the floor is taken
\emph{after} ramification scaling. The product $I_uN_u$ means the integer-ring
span of all $ay$ with $a\in I_u$, $y\in N_u$, exactly the source
\texttt{nPrime}.

Replace the source-lattice field of $D$ by $N'_u$, its height field by
the prescribed $h_u=\texttt{hExpl6}(D,L,Y.X,C_{\rm sz},u)$, and its
common height allowance by
\[
                  H'=\max\{1,\operatorname{sFrame}(N',h)
                                      +\mu_{\max}(RR)\}.       \tag{RF.21}
\]
The vector spaces, forms, host, weights, local comparisons and local
selected modules are retained. Hence the defining $N'$ and $h$ fields
now hold as identities. The numerical functions $\mathcal V_p,G,\tau,
\iota,TT_p$ and the local module $NP$ are unchanged, since none reads the
replaced output-lattice or height field. The auxiliary row-side maps are
unchanged as well: induction on their defining chain leaves each map,
its image and its adjoint composition identical. Thus every supplied
comparison, kernel and auxiliary datum transports to the new record.
This constructs the literal point \texttt{LedPart.pt} and preserves
its scale, degree identities, \texttt{Val} and \texttt{CanonInst}.
If the supplied recursive datum has the tame, Minkowski and local
quadratic-twisting properties, these transport too. Installing their
records does not prove their geometric or metric premises.

\paragraph{RF9: the two eventual scale inequalities have elementary proofs.}
For fixed $\epsilon>0$, \texttt{RankVal} gives
$n\leq2x^{2+8\epsilon}\leq x^3$ when $x\geq4$, because
$\epsilon<1/52$. Also
\[
                         w''=\frac{e^{x^\epsilon-x^{\epsilon/3}}}{n}.
\]
For $x^{2\epsilon/3}\geq2$ and $\log x\geq12/\epsilon^2$,
$x^\epsilon-x^{\epsilon/3}\geq x^\epsilon/2\geq3\log x$.
The last inequality follows from
$e^{\epsilon\log x}\geq(\epsilon\log x)^2/2$.
Thus $w''\geq x^3/n\geq1$, with an explicit threshold depending only
on $\epsilon$.
Next, if $c_s>0$, $m\geq1$, $T>0$, and
$\log T\geq c_smx^\epsilon$, then
\[
 x^{2\epsilon/3}\geq c_s^{-1}
       \quad\Longrightarrow\quad \log T\geq x^{\epsilon/3}
       \quad\Longrightarrow\quad X'\leq T.                   \tag{RF.22}
\]
These prove the two threshold lemmas used in the point join without
changing any previously fixed local constant.

\paragraph{RF10: the uniform point join and its exact remaining suppliers.}
The recursive premise \texttt{RecExP} has a useful order of choices.
After the external constants it supplies $c_s>0$ and a threshold.
For each sufficiently large admissible column datum it chooses positive
$d_u,D_*$ and lower bounds on the logarithms of the desired type-prime
centers. It then supplies $D,Y$, \texttt{Over}, \texttt{RecOK},
\texttt{LedOK}, and the twisting data \emph{for every} level record
with those $d_u,D_*$ and the prescribed $\Gamma$. Its requested property
$Q$ holds for every compatible completed \texttt{LedPart}.
The universal quantifier over the level record permits the level choice
to be made after $d_u,D_*$ have been fixed.

Retain the following three substantive suppliers with their literal
uniform hypotheses: \texttt{chGood\_of\_host}, which constructs the
compatible weight columns from the host and selected degrees;
\texttt{exists\_levCh\_of\_levelPrime}, which consumes
\texttt{LevelPrime} and chooses the prescribed norm/congruence primes;
and \texttt{exists\_sizeData\_over}, which constructs $L_0$ and the
same-host/word-label identifications from the recursive datum. These
are not consequences of the refill. With them, the point join is now
explicit. Enlarge the middle constant thresholds to include $c_g\geq10$,
replace the lower $c_D$ threshold by its maximum with $0$, and the lower
$c_h$ threshold by its maximum with $402$. Take the maximum of the
recursive, three supplier and two RF9 thresholds. For each admissible
host, selected degree pair, word code and distinguished position:
construct the columns; obtain $d_u,D_*$ recursively; choose the levels;
apply the universal recursive assertion; obtain $L_0$; apply RF7; and
install the point by RF8. The recursive property $Q$ applies to this
very \texttt{LedPart}, and the retained tame, Minkowski and twisting
properties give the accompanying duties.

The threshold order is exactly
\[
 \epsilon_0;\quad(\epsilon,C_8,C_5,c_6,C_b);\quad K_0;\quad C;
 \quad c_{D,0};\quad c_D;\quad c_{h,0};\quad c_h;\quad x_0;
 \quad(s,R,H,q_0,a_T,b_T,\mathrm{code},u_0).                    \tag{RF.23}
\]
No threshold may depend on the final degree pair or distinguished
position. The resulting point retains the original host field, top and
second sets, primes, Frobenius labels, word code and $u_0$ through the
supplied \texttt{Over}/initial-size identifications. This is the
\texttt{objCP\_and\_of\_recExP} join, with its three construction
suppliers still explicit. It does not assert \texttt{RecExP},
\texttt{LevelPrime}, or \texttt{NumAt}. The exact next-source list
records those remaining constructions rather than treating the record
type \texttt{ObjCP} as a proof of their existence.

\subsection{Actual numerical estimates beyond the refill}

These estimates have the precise \texttt{SzRow} order of quantifiers:
for each $0<\epsilon<1/52$ and fixed standard constant record $C$, choose
$c_{D,0}$; for every $c_D\geq c_{D,0}$ choose $c_{h,0}$; for every
$c_h\geq c_{h,0}$ choose $x_0$; then the conclusions hold for every
supplied $D,L$ with $\operatorname{Val}(C,c_D,c_h,L)$,
$L.\mathrm{sc}.\epsilon=\epsilon$, and $x\geq x_0$.
There is no choice of a different $D$ or $L$ in these assertions.
We write $\mathcal L=x^\epsilon$, $\delta=x^{-\epsilon/8}$,
$r=k+m-1$, $f=2n$, and $N_m=k+m(t+1)$.
All terminating decimals below denote exact rational numbers.

\paragraph{RF11: an explicit comparison principle for all the scales.}
If $a<b$, $p$ is a nonnegative integer and $A$ is fixed, then
\[
             A x^a(\log x)^p=o(x^b)\quad(x\longrightarrow\infty).
                                                               \tag{RF.24}
\]
Indeed, with $y=\log x\geq1$ and $d=b-a>0$, the positive series for
$e^{dy}$ gives
\[
 \frac{|A|y^p}{e^{dy}}
       \leq \frac{|A|(p+1)!}{d^{p+1}y}\longrightarrow0.
\]
This supplies an actual threshold for any requested positive upper
constant. Taking the maximum of three thresholds proves, for
$0\leq a<b$ and $\kappa>0$,
\[
 c\log x+A x^a(\log x)^p+A\leq\kappa x^b                 \tag{RF.25}
\]
eventually, with absolute values used for negative coefficients.
Consequently, for any fixed real $c$ and $\kappa>0$, all the following
hold at one threshold depending only on $\epsilon,c,\kappa$:
\begin{gather*}
 x^c\leq e^{x^{\epsilon/8}},\qquad
 x^c\leq (X')^\kappa,\qquad x^cX'\leq b^\kappa,
 \qquad b'x^c\leq b,\\
 b^{1-\kappa}x^c\leq b',\qquad
 b^{1-\kappa}x^c\leq w''\leq b'
       \quad(1\leq n\leq x^3),\\
 e^{c x^{\epsilon/5}(\log x)^2}x^c\leq(X')^\kappa.
                                                               \tag{RF.26}
\end{gather*}
For example the lower bound for $w''$ reduces, using
$\log n\leq3\log x$, to
$(c+3)\log x+x^{\epsilon/3}\leq\kappa x^\epsilon$.
The other inequalities follow by taking logarithms in (RF.1).
Thus this comparison principle proves the actual scale-polynomial
supplier used below, including its uniformity in $n$.

\paragraph{RF12: ranks, the large order, and the seven rank-parameter bounds.}
Put
\[
 T=\lceil x^{377}\rceil_+,
 \quad W=10r\left(f+\frac{f^2(k-1)^3\mathcal L^2}{m}\right),
 \quad n_W=8\lceil W/8\rceil_+ .                            \tag{RF.27}
\]
The supplied rank condition says $r+t=8fr+4+Tn_W$.
For $x\geq1000^{1000}$, constants at most $1000$ can be absorbed into
an additional exponent $0.001$: if $A\leq1000$ and $u+0.001\leq v$,
then $Ax^u\leq x^v$. The three defining rank windows therefore give
\[
\begin{array}{c|ccccccc}
\text{quantity}&f&r&k&\mathcal L&m&m^{-1}&f\binom r2\\ \hline
\text{upper power of }x&2.156&4.694&4.694&0.020&0.501&-0.498&11.544
\end{array}                                                   \tag{RF.28}
\]
For instance $r\leq2k\leq4x^{9/2+10\epsilon}$ and
$m\geq x^{1/2-\epsilon/16}\geq x^{0.498}$; also
$f\binom r2\leq fr^2$.
The second term of $W$ has exponent at most
$4.694+2(2.156)+3(4.694)+2(0.020)-0.498=22.630$.
The first has exponent at most $6.850$. Since $n_W<W+8$,
\[
 n_W\leq28x^{22.7}\leq x^{22.9},\quad
 x^{377}\leq T\leq x^{378},\quad
 x^{377}\leq t\leq13x^{400.9}\leq x^{401}.                  \tag{RF.29}
\]
The lower bound for $t$ uses $n_W\geq1$ and $8fr\geq r$.
Finally $N_m=k+m(t+1)\leq3x^{401.6}\leq x^{402}$.
This proves all seventeen conclusions of the source's
\texttt{rank\_basic\_bounds}: in addition to (RF.28)--(RF.29),
\[
\begin{gathered}
 f\leq x^3,\quad m\leq x,\quad r\leq x^5,\quad k\leq r\leq2k,
 \quad1\leq\mathcal L\leq x,\\
 x^{6.5}\leq fr\mathcal L\leq x^{6.87},\quad
 f\binom r2\leq x^{12},\quad kf/m\leq x^7,\\
 x^{1/2-\epsilon/16}\leq m\leq2x^{1/2-\epsilon/16}.
\end{gathered}                                                 \tag{RF.30}
\]
The lower bound for $fr\mathcal L$ follows from
$f\geq x^2$, $r\geq k\geq x^{9/2}$ and $\mathcal L\geq1$;
the upper bound and $kf/m$ follow by adding the exponents in (RF.28).

Let $R_0=nx^{-\epsilon}$ and let $C_k,C_6,C_f$ be the actual fixed
constants. At a further threshold the seven required inequalities are
\[
\begin{gathered}
 C_km^5f\mathcal L x^\epsilon\leq k,\qquad
 k\leq x^{1/2}nR_0,\qquad \sqrt{k/m}\leq R_0,\\
 m^2\mathcal L\leq k,\quad mx^{1/2}\log x\leq k,\quad
 C_6m^6\mathcal L\leq k\log k,\quad C_f\mathcal L\leq f\log k.
\end{gathered}                                                 \tag{RF.31}
\]
For the first three, the ratios of the left sides to the available
lower bounds on the right sides are bounded respectively by
$128|C_k|x^{-5\epsilon/16}$,
$2x^{-5\epsilon}$, and, after squaring the third inequality,
$2x^{-63\epsilon/16}$. They tend to zero.
For the last four put $y=x^{1/2}$ and require
$y\geq\max(8,64C_6,C_f)$. Then
$m\leq2y$, $k\geq y^9$, $f\geq y^4$,
$\mathcal L\leq y$, $0\leq\log x\leq y^2$, and $\log k\geq1$.
The four left sides are at most $4y^3$, $2y^4$, $y^8$, and $y^2$,
respectively. This proves the literal rank-parameter families, not only
their eventual conjunction.

\paragraph{RF13: thin-size domination and every prescribed column gap.}
From (RF.26) with $c=5|C_w|$ and $\kappa=1/10$, and $1\leq k\leq x^5$,
\[
                    k^{C_w}\leq w'',\qquad b^{9/10}\leq w''.
                                                               \tag{RF.32}
\]
Indeed both quantities are at most the product
$b^{9/10}x^{5|C_w|}$, which is at most $w''$.
This retains the source's fixed value $c_w=9/10$.

For the gap conditions let
$\Delta_{\mathrm{cl}}=\lceil x^{c_D}\rceil_+$,
$H_0=2\lceil x^{c_h}\rceil_+$ and
$G=C_{\mathrm{gap}}km^4+N_m m^{10}$.
Take $c_D,c_h\geq430$, $x\geq\max(2,C_{\mathrm{gap}})$,
and enlarge the scale threshold until $b,b',w''\geq x^{430}$.
By (RF.30),
$G\leq x^{10}+x^{412}<x^{430}$.
The actual \texttt{HCVal} columns are
\[
 v_i=(k-2-i)4\Delta_{\mathrm{cl}},\qquad b_j=b_1+jH_0,
       \qquad e_j=b_j+H_0/2.                                 \tag{RF.33}
\]
At a top, $b_1=b_T+\lfloor b'/x\rfloor\geq b'$; at a thick
position $b_1\geq\lceil b\rceil$; at a thin position
$b_1=\lceil w''\rceil$. Therefore $b_1\geq x^{430}$ in all cases.
The three gaps are exactly $4\Delta_{\mathrm{cl}}$, $v_{k-2}+b_1=b_1$, and $H_0$.
Each is strictly greater than $G$. This proves all three clauses of
\texttt{szRow\_sig\_gaps} with the actual columns.

\paragraph{RF14: the full band, lambda and host-prime bounds.}
Here the integer orbit-count construction is also expanded. Set
$h=2+8\epsilon$, $r_G=\lceil8(h+2)/\epsilon\rceil_++1$,
$q=6r_G+1$, and $\ell=\log x\geq1$. The source definitions are
\[
 U=\lceil1/\delta\rceil_+,\quad
 M=\lceil(h+3)10^q\ell\rceil_+,\quad
 G'=(M+5)qU,
 \quad N_1=\left\lceil\max(6eG'\delta,(h+3)\log_2x)\right\rceil_++1.
                                                               \tag{RF.34}
\]
Since $\delta U\leq2$,
$M+5\leq((h+3)10^q+6)\ell$, and $\log_2x\leq2\ell$, the
explicit constant
\[
 B_N=6e\bigl((h+3)10^q+6\bigr)2q+2(h+3)+2
\]
satisfies $1\leq N_1\leq B_N\log x$.
For $K_{\Lambda}\geq1$, put $B=K_{\Lambda}B_N^2$. The actual
$\Lambda=K_{\Lambda}N_1^2\log a$ consequently satisfies
\[
              x^\epsilon\leq\Lambda
                   \leq2B x^\epsilon(\log x)^2              \tag{RF.35}
\]
once $\log x\leq x^\epsilon$.

For any fixed real $c$, set
\[
 \rho(c)=\frac{8n\delta x^\epsilon\Lambda
                         +c\log x\,\Lambda}{k}.
\]
The rank windows, $\delta\leq1$ and (RF.35) bound its two numerators by
$32B x^{2+10\epsilon}(\log x)^2$ and
$2|c|B x^\epsilon(\log x)^3$.
By (RF.24), each is at most
$\tfrac12x^{5/2+10\epsilon}$ at a uniform threshold.
Dividing by $k\geq x^{9/2+10\epsilon}$ gives
\[
                         \rho(c)\leq x^{-2}.                \tag{RF.36}
\]
The choice $c=K_\rho K_N$ is exactly $\rho_B$; the choice
$c=C_\rho$ is exactly the parameter condition $\rho_P$.
Together with the supplied host top-count and $b_T$ bounds, this proves
the whole top/band/bottom parameter family as well.

Write $Q_c=QS\,q_0d$ for the host conductor product, using $q_0$ to
distinguish the host prime from the integer $q$ in (RF.34).
In this numerical argument $C_8=C.P.C_{p\mathrm{de}}$ is the constant
in the supplied \texttt{HostVal}.
The supplied \texttt{PairRow} at some $Y\leq x$ gives
\[
 \log Q_c\leq2x^{\epsilon/5}(\log x)^2,\qquad
 \log p_{\mathrm{de}},\ \log p'_{\mathrm{de}}
            \leq2C_8x^{\epsilon/5}(\log x)^2.
                                                               \tag{RF.37}
\]
Here the second statement follows from the actual inequalities
$p_{\mathrm{de}},p'_{\mathrm{de}}\leq Q_c^{C_8}$.
RF11 makes the last bound strictly smaller than $x^{\epsilon/3}$;
hence both auxiliary primes are strictly below $X'$.
For the already supplied level primes set $Y_g=(kb)^8$.
The host condition says $Y_g\leq g_i\leq2Y_g$. Since $k\leq x^5$,
RF11 gives
\[
 g_i\leq e^{9\mathcal L},\qquad \log Y_g\leq9\mathcal L,
 \qquad N_m^4Q_c^{C_Y}\leq Y_g^{1/4}.                       \tag{RF.38}
\]
For the last inequality take logarithms and use
$1608\log x+2C_Yx^{\epsilon/5}(\log x)^2\leq2x^\epsilon$;
the right logarithm is $2\log k+2x^\epsilon$.
These are the complete type-prime scale inequalities with the fixed
$C_{YL}=9$. They concern the supplied $g_i$ and do not choose any prime.

Finally, the lower and upper top-count bounds are precisely
\texttt{ChTT\_at}. The identity
$\#\mathrm{Word}=\lceil3\Gamma_L^2x^{C_F}\rceil_+$ implies
\texttt{ChCode\_at} with constants $(3,C_F)$ by the defining ceiling
inequality; it is a lower bound, not the separate exponential upper
bound on the word count. Since $1\leq J\leq C_8$,
$\epsilon/J\geq\epsilon/C_8$. Thus
$x\geq11^{5C_8/\epsilon}$ also gives
$11\leq\lfloor x^{(\epsilon/J)/5}\rfloor_+\leq\mathrm{splitBound}$.
This proves the full source \texttt{szRow\_band\_lambda\_primes}
and the additional split-bound clause, with no host existence inference.

\paragraph{RF15: construct every field of the actual block-length record.}
Suppose a compatible parameter record $P$ is supplied, so that
$P.C_{nW}=C.P.C_{nW}=10$. Set
\[
 \mathsf{TO}=T,\quad \mathsf n_W=n_W,\quad
 \mathsf n_P=8fr+4,\quad Z=fr\mathcal L,\quad
 \mathsf P=C_PZ^{12},\quad \mathsf D=f\binom r2,
 \quad s=r+6.                                                \tag{RF.39}
\]
Copy the seven constants from $C.T$. Standardness fixes
$e_P=12$, $C_p=1$, $c_{TO}=377$, gives $C_P\geq1000$, and gives the
other required lower bounds. Choose $c_D\geq800$, $c_h\geq1000$,
and then $x$ large enough that $C_P\leq x^{0.001}$,
$C_{hT},C_{bN}\leq x$ and RF11--RF14 apply.
Since $Z\leq x^{6.87}$ and $Z\geq x^{6.5}$,
\[
 1000x^{78}\leq\mathsf P\leq x^{82.441}\leq x^{85},
 \qquad \mathsf D\leq x^{11.544}\leq x^{12}.                \tag{RF.40}
\]
The defining second-set density gives
$\delta_2\geq x^{-\epsilon}$, hence
$\delta_2^{-1}\leq x^{0.020}$.
The strongest lower bound on $\mathsf{TO}$ now has an explicit exponent:
\[
 \frac{\mathsf P^4\mathsf D^3kf x}{m\delta_2}
 \leq x^{4(82.441)+3(11.544)+4.694+2.156+1-0.498+0.020}
 =x^{371.768}\leq x^{377}\leq\mathsf{TO}.                  \tag{RF.41}
\]
In particular this uses the exact order scale $377$; replacing the
inputs prematurely by much coarser bounds would not prove the claim.

The count interval is $W\leq\mathsf n_W<W+8$ and
$8\mid\mathsf n_W$. Since $m\leq k-1$, $f\geq2$ and
$\mathcal L\geq1$,
$W\geq10fr(1+(k-1)^2)\geq fr(r+1)$.
Also $(k-1)f+2r\leq2fr\leq W$ and $8fr+4\leq10fr\leq W$.
These prove all three lower-count conditions.
Moreover
$(2fm)^4\mathsf n_W\leq16x^{33.528}\leq\mathsf P$ and
$\mathsf n_W\leq\mathsf P$.
The exact splitting identity is the supplied rank identity in (RF.27).
Because $k,m$ are even, $r$ and $s=r+6$ are odd, and
$r+5<s\leq r+7$.

For completeness the remaining polynomial tests admit the following
explicit upper bounds. Increase $x_0$ to absorb the fixed coefficients.
\[
\begin{array}{l|r}
\text{quantity bounded above}&\text{power of }x\\ \hline
(\mathsf D+2)\mathsf n_W+2r&36\\
40(\mathsf D+2)&13\\
11\mathsf D^2&25\\
r+2(2\cdot5^4r)(\mathsf D+1)&18\\
\mathsf D k/m&17
\end{array}
\qquad
\begin{array}{l|r}
\text{quantity bounded above}&\text{power of }x\\ \hline
\mathsf{TO}^2&756\\
(\mathsf{TO}+\mathsf D+2)
 (\mathsf{TO}+\mathsf D+2+2N_m)&782\\
4\mathsf{TO}^2+2t&757\\
\mathsf P\mathsf D f m^2N_m\mathsf{TO}&882\\
\mathsf P(fmk)^3k/m&117\\
C_{hT}\mathsf{TO}N_m f^2m^2k\mathsf D\mathsf P&891
\end{array}                                                   \tag{RF.42}
\]
Every entry of the left table is at most $\mathsf{TO}$.
The first two entries of the right table are at most
$\Delta_{\mathrm{cl}}=\lceil x^{c_D}\rceil_+$; the last four are at most
$h_0=L.\etaSz/m=2\lceil x^{c_h}\rceil_+$.
For the product of two sums, use
$\mathsf{TO}+\mathsf D+2\leq x^{379}$ and the second factor
$\leq x^{403}$. For the first height entry use
$4x^{756}+2x^{401}\leq6x^{756}\leq x^{757}$.
The other height exponents are obtained by multiplying the explicit
bounds $\mathsf P\leq x^{85}$, $\mathsf D\leq x^{12}$,
$f\leq x^3$, $m\leq x$, $k\leq x^5$, $N_m\leq x^{402}$,
$\mathsf{TO}\leq x^{378}$; in the cubic entry use $m\geq1$.

The tests against $b$ and $\sqrt b$ are not additional assumptions:
\[
\begin{array}{l|r}
\mathsf{TO}\mathsf P x^{C_b}&x^{463+C_b}\\
\mathsf{TO}\mathsf P&x^{463}\\
\mathsf P^3\mathsf{TO}^2x^{C_b}&x^{1011+C_b}\\
\mathsf P^3\mathsf{TO}^2&x^{1011}\\
\mathsf{TO}\mathsf D^2\log b\,x^{C_b}&x^{403+C_b}\\
\mathsf{TO}\mathsf D^2\log b&x^{403}\\
C_{bN}rmk\mathsf P^2\mathsf D^3&x^{218}
\end{array}                                                   \tag{RF.43}
\]
RF11 puts the first, third, fifth and seventh rows below $b$, and the
second, fourth and sixth below $b^{1/2}$.
The three polynomial-size bounds are immediate from
$\mathsf P\leq x^{85}$, $\mathsf D\leq x^{12}$,
$\mathsf{TO}\leq x^{378}\leq x^{401}$.
This checks every defining field of \texttt{ParamsT}, including its
constant list and the identity $e_P=C.T.e_P$.
It proves the entire \texttt{szRow\_blockLengthParams} implication,
uniformly for every supplied compatible $P$.

\paragraph{The exact numerical supplier boundary.}
Including RF16 below, seven of the twenty-one parameter families in
\texttt{explicitConstants\_parameterConditions\_large} have now been
proved here: the three rank-$k$ bounds; the four other rank bounds;
top-count/band/bottom bounds; thin-size bounds; all signature gaps;
the type-prime scale bounds; and the upper word-cardinality bound.
RF15 proves the full block-length
supplier, and RF14 proves the full band/lambda/prime supplier.
The fourteen other parameter families and the complete
\texttt{SizesAt} package are still retained, with their exact source
targets recorded in the continuation files.

If those remaining suppliers give a compatible $P$ and
\texttt{SizesAt} for the same $D,L$, RF15 gives \texttt{ParamsT} and
RF14 gives every nontrivial remaining inequality in \texttt{NumAt}.
Its other fields are the literal identities in \texttt{Val}:
the constants, $b_T$, the top ideal $cc_K$, the logarithmic depth,
the common cloud vector, the rank windows, and the formula for $L_w$
using $C_t=401$. Thus \texttt{NumAt}, \texttt{ChTT\_at} and
\texttt{ChCode\_at} follow conditionally on precisely these two
remaining packages. Neither those packages nor the simultaneous
geometric construction is asserted by this numerical extension.

\paragraph{RF16: the literal upper bound on the number of words.}
This last estimate uses the actual exponential value of $\Gamma_L$.
Write $G=C.C_\Gamma\geq1$, $F=C.P.C_F\geq1$, and
$H=m^5f x^\epsilon$. The exact \texttt{Val} identities give
\[
 \Gamma_L=e^{GH},\qquad
 \#\mathrm{Word}=\left\lceil3e^{2GH}x^F\right\rceil_+.
                                                               \tag{RF.44}
\]
The rank windows, with the full saving in the exponent of $m$, imply
\[
 H\leq128x^{9/2+139\epsilon/16},\qquad
 k\geq x^{9/2+160\epsilon/16}.                              \tag{RF.45}
\]
The difference is $21\epsilon/16>0$. Fix a threshold at which
$k\geq6$,
\[
 256Gx^{9/2+139\epsilon/16}
                  \leq\tfrac18x^{9/2+160\epsilon/16},
 \qquad F\log x+\log6
                  \leq\tfrac18x^{9/2+160\epsilon/16}.        \tag{RF.46}
\]
For the first inequality it suffices that
$x^{21\epsilon/16}\geq2048G$; RF11 supplies the second.
Both thresholds depend only on $\epsilon,C$.
Put $A=3e^{2GH}x^F\geq3$. Since
$\lceil A\rceil_+<A+1\leq2A$,
\[
 \log\#\mathrm{Word}
       \leq\log6+2GH+F\log x\leq k/4
       \leq(k-3)\log2.                                     \tag{RF.47}
\]
The last inequality uses $k-3\geq k/2$ and $\log2>1/2$.
Exponentiating, and identifying the natural subtraction with the
integer $k-3$ because $k\geq6$, proves
\[
                         \#\mathrm{Word}\leq2^{k-3}.         \tag{RF.48}
\]
This is the complete \texttt{szRow\_card\_words\_le} family.
It has $c_{D,0}=c_{h,0}=0$, so it combines with the preceding
families without changing their quantifier order. The number of words
is not treated as polynomial: its logarithm is compared to the
larger rank $k$ with the explicit positive gap in (RF.45).
